
\documentclass{article}
\usepackage{colortbl}
\usepackage{makecell}
\usepackage{multirow}
\usepackage{supertabular}

\begin{document}

\newcounter{utterance}

\centering \large Interaction Transcript for game `taboo', experiment `medium\_en', episode 2 with reflexion{-}Qwen3.5{-}4B.
\vspace{24pt}

{ \footnotesize  \setcounter{utterance}{1}
\setlength{\tabcolsep}{0pt}
\begin{supertabular}{c@{$\;$}|p{.15\linewidth}@{}p{.15\linewidth}p{.15\linewidth}p{.15\linewidth}p{.15\linewidth}p{.15\linewidth}}
    \# & $\;$A & \multicolumn{4}{c}{Game Master} & $\;\:$B\\
    \hline

    \theutterance \stepcounter{utterance}  
    & & \multicolumn{4}{p{0.6\linewidth}}{
        \cellcolor[rgb]{0.9,0.9,0.9}{
            \makecell[{{p{\linewidth}}}]{
                \texttt{\tiny{[P1$\langle$GM]}}
                \texttt{You have attempted this same task before and did not succeed. These are the plans you wrote after those attempts.} \\
\texttt{Trial \#1: New plan: I failed because my clue was too specific and directly defined the target word ("A person who is expected to do something in the future"), allowing the opponent to guess it immediately. To win, I must provide a clue that hints at the concept without naming the role or using the related words (foresight, potential, future). Next time, I will describe the *action* or *state* associated with the word rather than the definition. For example, I will say "Someone who is about to become a doctor or employee" or "A future possibility that is likely to happen," ensuring the opponent needs to infer the specific noun rather than recognizing the definition.} \\
\\ 
\\ 
\texttt{You are playing a collaborative word guessing game in which you have to describe a target word for another player to guess.} \\
\\ 
\texttt{Rules:} \\
\texttt{(a) You have to reply in the form: CLUE: <some text>. Guesses from the other player will start with GUESS.} \\
\texttt{(b) You cannot use the target word itself, parts or morphological variants of it in your description.} \\
\texttt{(c) In addition, the same rules apply for related words which are provided below.} \\
\\ 
\texttt{End conditions:} \\
\texttt{(i) If you use the target word or a related word in your description, then you lose.} \\
\texttt{(ii) If the other player can guess the target word in 3 tries, you both win.} \\
\\ 
\texttt{Let us start.} \\
\\ 
\texttt{This is the target word that you need to describe and that the other player needs to guess:} \\
\\ 
\texttt{prospective} \\
\\ 
\texttt{Related words are:} \\
\\ 
\texttt{{-} foresight} \\
\texttt{{-} potential} \\
\texttt{{-} future} \\
\\ 
\texttt{Important: You are under time pressure, give short descriptions that are to the point!} \\
            }
        }
    }
    & \\ \\

    \theutterance \stepcounter{utterance}  
    & \multicolumn{4}{p{0.6\linewidth}}{
        \cellcolor[rgb]{0.9,0.9,0.9}{
            \makecell[{{p{\linewidth}}}]{
                \texttt{\tiny{[P1$\rangle$GM]}}
                \texttt{CLUE: A person who is about to start a job or study at a specific place.} \\
            }
        }
    }
    & & \\ \\

    \theutterance \stepcounter{utterance}  
    & & & \multicolumn{2}{p{0.3\linewidth}}{
        \cellcolor[rgb]{0.9,0.9,0.9}{
            \makecell[{{p{\linewidth}}}]{
                \texttt{\tiny{[GM$|$GM]}}
                \texttt{continue} \\
            }
        }
    }
    & & \\ \\

    \theutterance \stepcounter{utterance}  
    & & & \multicolumn{2}{p{0.3\linewidth}}{
        \cellcolor[rgb]{0.9,0.9,0.9}{
            \makecell[{{p{\linewidth}}}]{
                \texttt{\tiny{[GM$|$GM]}}
                \texttt{a person who is about to start a job or study at a specific place} \\
            }
        }
    }
    & & \\ \\

    \theutterance \stepcounter{utterance}  
    & & \multicolumn{4}{p{0.6\linewidth}}{
        \cellcolor[rgb]{0.9,0.9,0.9}{
            \makecell[{{p{\linewidth}}}]{
                \texttt{\tiny{[GM$\rangle$P2]}}
                \texttt{You are playing a collaborative word guessing game in which you have to guess a target word that another player describes to you.} \\
\\ 
\texttt{You can make one guess at each trial. You win when you guess the target word. You lose when you cannot guess it in 3 tries.} \\
\\ 
\texttt{After each trial you will get a new hint from the other player which starts with CLUE.} \\
\\ 
\texttt{Make your guesses by just saying the word using the following form: GUESS: <a word>} \\
\\ 
\texttt{Let us start.} \\
\\ 
\\ 
\texttt{CLUE: A person who is about to start a job or study at a specific place.} \\
            }
        }
    }
    & \\ \\

    \theutterance \stepcounter{utterance}  
    & & & \multicolumn{4}{p{0.6\linewidth}}{
        \cellcolor[rgb]{0.9,0.9,0.9}{
            \makecell[{{p{\linewidth}}}]{
                \texttt{\tiny{[GM$\langle$P2]}}
                \texttt{GUESS: applicant} \\
            }
        }
    }
     \\ \\

    \theutterance \stepcounter{utterance}  
    & & & \multicolumn{2}{p{0.3\linewidth}}{
        \cellcolor[rgb]{0.9,0.9,0.9}{
            \makecell[{{p{\linewidth}}}]{
                \texttt{\tiny{[GM$|$GM]}}
                \texttt{continue} \\
            }
        }
    }
    & & \\ \\

    \theutterance \stepcounter{utterance}  
    & & & \multicolumn{2}{p{0.3\linewidth}}{
        \cellcolor[rgb]{0.9,0.9,0.9}{
            \makecell[{{p{\linewidth}}}]{
                \texttt{\tiny{[GM$|$GM]}}
                \texttt{applicant} \\
            }
        }
    }
    & & \\ \\

    \theutterance \stepcounter{utterance}  
    & & \multicolumn{4}{p{0.6\linewidth}}{
        \cellcolor[rgb]{0.9,0.9,0.9}{
            \makecell[{{p{\linewidth}}}]{
                \texttt{\tiny{[P1$\langle$GM]}}
                \texttt{GUESS: applicant} \\
            }
        }
    }
    & \\ \\

    \theutterance \stepcounter{utterance}  
    & \multicolumn{4}{p{0.6\linewidth}}{
        \cellcolor[rgb]{0.9,0.9,0.9}{
            \makecell[{{p{\linewidth}}}]{
                \texttt{\tiny{[P1$\rangle$GM]}}
                \texttt{CLUE: Someone who is expected to become a doctor, lawyer, or employee in the near future.} \\
            }
        }
    }
    & & \\ \\

    \theutterance \stepcounter{utterance}  
    & & & \multicolumn{2}{p{0.3\linewidth}}{
        \cellcolor[rgb]{0.9,0.9,0.9}{
            \makecell[{{p{\linewidth}}}]{
                \texttt{\tiny{[GM$|$GM]}}
                \texttt{continue} \\
            }
        }
    }
    & & \\ \\

    \theutterance \stepcounter{utterance}  
    & & & \multicolumn{2}{p{0.3\linewidth}}{
        \cellcolor[rgb]{0.9,0.9,0.9}{
            \makecell[{{p{\linewidth}}}]{
                \texttt{\tiny{[GM$|$GM]}}
                \texttt{Related word 'future' (stem=futur) is similar to clue word 'future' (stem=futur)} \\
            }
        }
    }
    & & \\ \\

    \theutterance \stepcounter{utterance}  
    & \multicolumn{4}{p{0.6\linewidth}}{
        \cellcolor[rgb]{0.95,0.95,0.95}{
            \makecell[{{p{\linewidth}}}]{
                \texttt{\tiny{[P2$\rightarrow$P1]}}
                \texttt{Post{-}episode Reflexion (not sent to GM):} \\
\texttt{New plan: I failed because my clues were still too close to the definition of the target word ("prospective"), allowing the opponent to deduce the specific noun immediately. To win, I must describe the *abstract concept* or *metaphorical meaning* of the word without mentioning the person, the profession, or the related words (foresight, potential, future). Next time, I will focus on the idea of "something that is likely to happen but not yet real" or "a future possibility," forcing the opponent to bridge the gap between the abstract concept and the specific noun "prospective."} \\
            }
        }
    }
    & & \\ \\

    \theutterance \stepcounter{utterance}  
    & & & \multicolumn{4}{p{0.6\linewidth}}{
        \cellcolor[rgb]{0.95,0.95,0.95}{
            \makecell[{{p{\linewidth}}}]{
                \texttt{\tiny{[P1$\rightarrow$P2]}}
                \texttt{Post{-}episode Reflexion (not sent to GM):} \\
\texttt{New plan: I guessed "applicant" immediately without considering that the clue "A person who is about to start a job or study at a specific place" could also refer to a "student" or "intern," and I only had one guess left after that. Next time, I will analyze the clue to generate multiple potential candidates (e.g., student, intern, recruit, cadet) and select the most common or specific one for my first guess, ensuring I use a different word for the second and third guesses if the first fails.} \\
            }
        }
    }
     \\ \\

\end{supertabular}
}

\end{document}
